\begin{ouverture}[courbe]{1}{Calculer dans les complexes}{Nombres complexes}{Forme algébrique · Conjugaison · Binôme}
\objectif{Calculer avec $\mathrm i^2=-1$.}
\objectif{Utiliser la conjugaison et l’inverse.}
\objectif{Résoudre des équations avec $z$ et son conjugué.}
\objectif{Appliquer et démontrer la formule du binôme.}
\prerequis{Calcul littéral, coefficients binomiaux et systèmes de deux équations.}
\end{ouverture}
\essentiel

\essSection{1}{Un nombre à deux coordonnées}
\begin{definition}
Tout nombre complexe s'écrit de manière unique $z=x+\mathrm iy$, avec $x,y\in\R$ et $\mathrm i^2=-1$.
$x=\Re(z)$ est la partie réelle et $y=\Im(z)$ la partie imaginaire : toutes deux sont des nombres réels.
Deux complexes sont égaux si leurs parties réelles et imaginaires sont égales.
\end{definition}
Les réels s'identifient aux complexes de partie imaginaire nulle. Les imaginaires purs ont une partie réelle nulle, y compris zéro.
L'addition et le produit suivent les règles usuelles de distributivité avec $\mathrm i^2=-1$.
\begin{exemple}
$(2+3\mathrm i)(1-2\mathrm i)=2-4\mathrm i+3\mathrm i-6\mathrm i^2=8-\mathrm i$.
Les puissances de $\mathrm i$ se répètent par période $4$ : $1,\mathrm i,-1,-\mathrm i$.
\end{exemple}
\essSection{2}{La conjugaison}
\begin{definition}
Le conjugué de $z=x+\mathrm iy$ est $\bar z=x-\mathrm iy$.
On a $\overline{\bar z}=z$, $z+\bar z=2\Re(z)$ et $z-\bar z=2\mathrm i\Im(z)$.
\end{definition}
\begin{propriete}
$\overline{z+w}=\bar z+\bar w$, $\overline{zw}=\bar z\bar w$ et $\overline{z^n}=\bar z^{\,n}$ pour $n\in\N$.
Un complexe est réel si et seulement si $z=\bar z$ ; il est imaginaire pur si et seulement si $z=-\bar z$.
\end{propriete}
\begin{demonstration}
Pour $z=x+\mathrm iy$, $w=u+\mathrm iv$, le produit vaut $(xu-yv)+\mathrm i(xv+yu)$. Son conjugué est précisément $(x-\mathrm iy)(u-\mathrm iv)$. La règle des puissances suit par récurrence sur $n$, à partir de $n=0$.
\end{demonstration}
\figurelivre{e01-conjugaison.pdf}{Dans le plan complexe, $z=2+3\mathrm i$ et $\bar z=2-3\mathrm i$ sont symétriques par rapport à l'axe réel. L'ordonnée représente $\Im(z)$, pas $\mathrm i\Im(z)$.}
\essSection{3}{Diviser et résoudre une équation}
Pour $z=x+\mathrm iy\ne0$, $z\bar z=x^2+y^2>0$, donc
$1/z=\bar z/(x^2+y^2)$.
Pour mettre un quotient sous forme algébrique, multiplier le numérateur et le dénominateur par le conjugué du dénominateur.
\begin{exemple}
$\dfrac{1+\mathrm i}{2-\mathrm i}=\dfrac{(1+\mathrm i)(2+\mathrm i)}5=\dfrac15+\dfrac35\mathrm i$.
Pour résoudre $z+2\bar z=3+4\mathrm i$, poser $z=x+\mathrm iy$ : $3x-\mathrm iy=3+4\mathrm i$, donc $x=1,y=-4$.
\end{exemple}
L'équation $az=b$ a une unique solution $b/a$ si $a\ne0$. Si $a=0$, elle a toutes les solutions lorsque $b=0$, et aucune sinon.
Conjuguer $z(1/z)=1$ donne $\bar z\,\overline{1/z}=1$, donc $\overline{1/z}=1/\bar z$.
\essSection{4}{Le binôme reste valable}
\begin{theoreme}
Pour $a,b\in\C$ et $n\in\N$,
$\displaystyle(a+b)^n=\sum_{k=0}^n\binom nk a^{n-k}b^k$.
\end{theoreme}
\begin{demonstration}
L'identité vaut pour $n=0$. Supposer l'identité au rang $n$, multiplier par $a+b$ et regrouper les coefficients d'une même puissance de $b$. L'identité de Pascal $\binom nk+\binom n{k-1}=\binom{n+1}k$ donne la formule au rang $n+1$. Les facteurs complexes commutent, donc le regroupement est valide.
\end{demonstration}
\begin{attention}
Il n'existe pas d'ordre sur $\C$ permettant de comparer tous les complexes comme des réels tout en conservant les règles usuelles des inégalités. On ne résout pas « $z>0$ » pour un complexe arbitraire. Les conditions de réalité se traduisent en revanche par des équations sur $x,y$.
\end{attention}

\exercices
\markright{Prendre confiance}
\begingroup\renewcommand{\otBesoinPlace}[1]{}\begin{minipage}{\linewidth}
\rubrique{Prendre confiance}
\begin{exercice}[1]{Parties réelle et imaginaire}

\begin{questions}
\item Donner les parties réelle et imaginaire de $3-5\mathrm i$.
\item Calculer $(2+\mathrm i)+(4-3\mathrm i)$.
\item Résoudre $x+2\mathrm iy=3-4\mathrm i$ avec $x,y$ réels.
\end{questions}
\end{exercice}
\begin{exercice}[1]{Produits}

\begin{questions}
\item Calculer $(1+\mathrm i)^2$.
\item Calculer $(2+3\mathrm i)(1-2\mathrm i)$.
\item Calculer $\mathrm i^{2026}$.
\end{questions}
\end{exercice}
\begin{exercice}[1]{Conjuguer}

\begin{questions}
\item Conjuguer $-2+7\mathrm i$.
\item Calculer $z+\bar z$ pour $z=4-3\mathrm i$.
\item Calculer $z\bar z$ pour ce même $z$.
\end{questions}
\end{exercice}
\begin{exercice}[1]{Inverses}

\begin{questions}
\item Mettre $1/(1+\mathrm i)$ sous forme algébrique.
\item Faire de même pour $1/(3-4\mathrm i)$.
\item Vérifier l’un des résultats par multiplication.
\end{questions}
\end{exercice}
\end{minipage}\endgroup
\clearpage
\markright{Appliquer et relier}
\begingroup\renewcommand{\otBesoinPlace}[1]{}\begin{minipage}{\linewidth}
\rubrique{Appliquer et relier}
\begin{exercice}[2]{Équation linéaire}

\begin{questions}
\item Résoudre $(1+\mathrm i)z=2$.
\item Résoudre $2z-3=\mathrm i$.
\item Discuter $0z=b$ suivant $b$.
\end{questions}
\end{exercice}
\begin{exercice}[2]{Un quotient}

\begin{questions}
\item Calculer $(3+\mathrm i)/(1-\mathrm i)$.
\item Vérifier en remultipliant par $1-\mathrm i$.
\item En déduire le conjugué du quotient.
\end{questions}
\end{exercice}
\begin{exercice}[2]{Avec le conjugué}

\begin{questions}
\item Résoudre $z+2\bar z=3+4\mathrm i$.
\item Résoudre $z-\bar z=6\mathrm i$.
\item Pourquoi la seconde équation a-t-elle une infinité de solutions ?
\end{questions}
\end{exercice}
\begin{exercice}[2]{Binôme}

\begin{questions}
\item Développer $(1+\mathrm i)^3$ par le binôme.
\item Calculer $(1+\mathrm i)^4$.
\item En déduire $(1+\mathrm i)^8$.
\end{questions}
\end{exercice}
\end{minipage}\endgroup
\clearpage
\markright{Chercher et argumenter}
\begingroup\renewcommand{\otBesoinPlace}[1]{}\begin{minipage}{\linewidth}
\rubrique{Chercher et argumenter}
\begin{exercice}[3]{Réalité d’un quotient}
$z=(x+\mathrm i)/(1+\mathrm i)$ avec $x\in\R$.
\begin{questions}
\item Mettre $z$ sous forme algébrique.
\item Pour quelle valeur de $x$ est-il réel ?
\item Pour quelle valeur est-il imaginaire pur ?
\end{questions}
\end{exercice}
\begin{exercice}[3]{Conjugaison et puissance}

\begin{questions}
\item Prouver $\overline{z^2}=\bar z^2$.
\item Étendre la propriété par récurrence à tout entier naturel.
\item L’étendre aux puissances entières négatives lorsque $z\ne0$.
\end{questions}
\end{exercice}
\begin{exercice}[3]{Une équation à analyser}
$z\bar z=2z$.
\begin{questions}
\item Poser $z=x+\mathrm iy$ et comparer les parties.
\item Résoudre.
\item Vérifier les solutions dans l’équation initiale.
\end{questions}
\end{exercice}
\begin{exercice}[3]{Une somme binomiale}

\begin{questions}
\item Exprimer $(1+\mathrm i)^6$ avec les coefficients binomiaux.
\item Calculer sa forme algébrique par une méthode plus courte.
\item En déduire la somme des termes réels du développement.
\end{questions}
\end{exercice}
\end{minipage}\endgroup
